% ===============================================================
% APRESENTAÇÃO (BEAMER) - TCC2 - SENAC SANTO AMARO
% ===============================================================
% Defesa final do Trabalho de Conclusão de Curso (TCC2), com o
% conteúdo e os números já fechados em tcc_2_adapted.tex.
%
% >>> DURAÇÃO-ALVO DA APRESENTAÇÃO: ~20 MINUTOS <<<
%
% Por isso, em relação ao documento escrito, foram DELIBERADAMENTE
% OMITIDOS da fala principal: o referencial teórico em profundidade
% (AST, CFG, DFA, arquitetura Transformer, NLP, agentes), o
% detalhamento de cada trabalho correlato, o detalhamento fino dos
% formatos de regra/combinações de modelo (fica só um resumo rápido,
% o nível de detalhe completo está no apêndice) e o cronograma (já
% cumprido). Esse material fica no APÊNDICE, como apoio para a
% arguição da banca (não conta no tempo).
%
% Duas ferramentas centrais (SemGrep, Juliet Test Suite) e o
% vocabulário das configurações (formato/combo/example_mode/
% guard/project_apis) FORAM TRAZIDOS de volta para a fala principal,
% a pedido explícito do autor -- sem eles, os rótulos das tabelas de
% resultado ("Combo D, template, pairs+guard") não fazem sentido para
% quem não leu o TCC. O vocabulário de configurações foi desmembrado
% em um slide por modo de exemplos (findings/pairs/graph) a pedido do
% autor, em vez de um único slide resumido.
%
% IMAGENS: três gráficos (CWE-416, CWE-415, comparação com a regra
% tradicional) ainda não foram inseridos -- os slides correspondentes
% têm um AVISO "[Espaço reservado para gráfico -- Images/...svg]" no
% lugar da imagem. Troque pelo \includegraphics correspondente quando
% os arquivos estiverem prontos (os .svg já existem em Images/, mas
% precisam ser convertidos para PDF/PNG antes de entrar aqui).
%
% LOGO: a logo do SENAC aparece GRANDE no primeiro slide (capa) e
% PEQUENA no canto superior direito dos demais slides.
% Arquivo já encontrado em: assets/logo-senac.png
%
% COMO COMPILAR:
%   pdflatex apresentacao_tcc2.tex
%   pdflatex apresentacao_tcc2.tex   (2x para fixar índice/numeração)
% ===============================================================

\documentclass[
    aspectratio=169,
    12pt,
    brazil
]{beamer}

% ---------------------------------------------------------------
% PACOTES
% ---------------------------------------------------------------
\usepackage[T1]{fontenc}
\usepackage[utf8]{inputenc}
\usepackage[brazil]{babel}
\usepackage{graphicx}
\usepackage{booktabs}
\usepackage{multirow}
\usepackage{tikz}
\usetikzlibrary{positioning}
\usepackage{pgfgantt}
\usepackage{appendixnumberbeamer}

% ---------------------------------------------------------------
% TEMA E CORES
% ---------------------------------------------------------------
\usetheme{default}
\useinnertheme{circles}
\setbeamertemplate{navigation symbols}{}

\definecolor{senacblue}{RGB}{0,51,102}
\definecolor{senacorange}{RGB}{243,146,0}

\setbeamercolor{structure}{fg=senacblue}
\setbeamercolor{frametitle}{fg=white,bg=senacblue}
\setbeamercolor{title}{fg=senacblue}
\setbeamercolor{section in toc}{fg=senacblue}
\setbeamercolor{itemize item}{fg=senacorange}
\setbeamercolor{itemize subitem}{fg=senacorange}
\setbeamerfont{frametitle}{series=\bfseries}
\setbeamertemplate{frametitle}[default][left]

\setbeamertemplate{footline}{%
  \leavevmode%
  \hbox{%
    \begin{beamercolorbox}[wd=.34\paperwidth,ht=2.4ex,dp=1ex,leftskip=1em]{author in head/foot}%
      \usebeamerfont{author in head/foot}\insertshortauthor
    \end{beamercolorbox}%
    \begin{beamercolorbox}[wd=.5\paperwidth,ht=2.4ex,dp=1ex,center]{title in head/foot}%
      \usebeamerfont{title in head/foot}\insertshorttitle
    \end{beamercolorbox}%
    \begin{beamercolorbox}[wd=.16\paperwidth,ht=2.4ex,dp=1ex,rightskip=1em]{title in head/foot}%
      \usebeamerfont{title in head/foot}\hfill\insertframenumber/\inserttotalframenumber
    \end{beamercolorbox}%
  }\vskip0pt%
}
\setbeamercolor{author in head/foot}{bg=senacblue,fg=white}
\setbeamercolor{title in head/foot}{bg=senacblue!85,fg=white}

% ---------------------------------------------------------------
% LOGO DO SENAC (grande na capa, pequena no canto dos demais slides)
% ---------------------------------------------------------------
\newcommand{\logosenacgrande}{%
  \IfFileExists{assets/logo-senac.png}{%
    \includegraphics[height=2.6cm]{assets/logo-senac.png}%
  }{%
    \begin{tikzpicture}
      \node[draw=senacblue,line width=1pt,rounded corners,
            minimum width=4.5cm,minimum height=2.4cm,align=center,
            fill=senacblue!5]
        {\bfseries\color{senacblue}\Large SENAC\\[2pt]
         \normalsize Santo Amaro\\[6pt]
         \scriptsize\itshape (coloque assets/logo-senac.png)};
    \end{tikzpicture}%
  }%
}

\newcommand{\logosenacpequena}{%
  \IfFileExists{assets/logo-senac.png}{%
    \includegraphics[height=0.8cm]{assets/logo-senac.png}%
  }{%
    \fbox{\scriptsize\bfseries\color{senacblue}SENAC}%
  }%
}

% Logo pequena no canto superior direito dos slides com título (chip
% branco para a logo preta ficar visível sobre a barra azul). Frames
% "plain" (capa/encerramento) não têm frametitle e ficam sem a logo.
\addtobeamertemplate{frametitle}{}{%
  \begin{tikzpicture}[remember picture,overlay]
    \node[anchor=north east,fill=white,rounded corners=2pt,inner sep=3pt]
      at ([xshift=-6pt,yshift=-5pt]current page.north east) {\logosenacpequena};
  \end{tikzpicture}%
}

% ---------------------------------------------------------------
% DADOS DA APRESENTAÇÃO (preencha)
% ---------------------------------------------------------------
\title[Síntese de regras SemGrep com LLMs]{Análise Estática Generativa com LLMs: Otimizando a Detecção de Vulnerabilidades no SemGrep via Abordagem Multi-Agente}
\subtitle{Apresentação final --- TCC2}
\author[Caio Xavier da Silva]{Caio Xavier da Silva}
\institute[SENAC Santo Amaro]{%
  SENAC --- Unidade Santo Amaro\\
  Bacharelado em Ciência da Computação\\
  \footnotesize Orientador: Prof. Eduardo Heredia \\ Coorientador: Prof. Thyago Conchado Quintas}
\date{São Paulo, 2026}

% ===============================================================
\begin{document}

% ---------------------------------------------------------------
% SLIDE 1 - CAPA (logo GRANDE)
% ---------------------------------------------------------------
{
\setbeamertemplate{footline}{}
\begin{frame}[plain]
  \centering
  \vspace{0.1cm}
  \logosenacgrande

  \vspace{0.2cm}
  {\Large\usebeamercolor[fg]{title}\bfseries\inserttitle\par}
  \vspace{0.1cm}
  {\normalsize\insertsubtitle\par}

  \vspace{0.3cm}
  {\insertauthor\par}
  \vspace{0.1cm}
  {\footnotesize\insertinstitute\par}

  \vfill
  {\footnotesize\insertdate\par}
\end{frame}
}

% ---------------------------------------------------------------
% SLIDE 2 - AGENDA
% ---------------------------------------------------------------
\begin{frame}{Roteiro da apresentação}
  \tableofcontents
\end{frame}

% ===============================================================
% Orçamento de tempo (alvo: ~20 minutos)
%   - Problema e motivação .......... ~3 min   (2 slides)
%   - Ferramentas utilizadas ......... ~2,5 min (3 slides: SemGrep, exemplo
%                                                 de regra, Juliet)
%   - Objetivos e método ............ ~2,5 min (2 slides)
%   - Resultados ..................... ~11,5 min (12 slides, incl. 7 slides
%                                                 de vocabulário: formato+
%                                                 combo, nossas combinações,
%                                                 por que detector pequeno/
%                                                 fixo, findings, pairs,
%                                                 graph, guard+project_apis;
%                                                 síntese por CWE inclui
%                                                 nota sobre recall vir só
%                                                 do teste isolado do Juliet)
%   - Conclusão ...................... ~1,5 min (1 slide)
% Capa + roteiro + encerramento não entram nesse orçamento.
% Referencial teórico em profundidade, trabalhos correlatos em
% detalhe, o detalhamento fino de formatos de regra/combos de modelo
% e o cronograma ficam no apêndice (sem tempo fixo -- usar apenas se
% a banca perguntar).
% ===============================================================

% ---------------------------------------------------------------
\section{Problema e motivação}
% ---------------------------------------------------------------
\begin{frame}{O problema}
  \begin{itemize}
    \item Uma falta de checagem de tamanho no OpenSSL causou o \textit{Heartbleed} (CVE-2014-0160).
    \item Entre 24\% e 55\% dos servidores HTTPS mais acessados do mundo ficaram vulneráveis, como: Yahoo, Stack Overflow e a Receita Federal do Canadá.
    \item Análise estática tradicional depende de regras rígidas e manuais -- menos de 5\% das lógicas de detecção usadas na indústria são escritas por usuários comuns.
    \item Submeter o código inteiro a um LLM: dezenas de milhões de \textit{tokens}, dias de processamento e risco de alucinação.
  \end{itemize}
\end{frame}

\begin{frame}{Pergunta de pesquisa}
  \begin{quote}
    \itshape Um \textit{framework} multi-agente baseado em LLMs é capaz de sintetizar regras para o SemGrep que detectem vulnerabilidades em código C com \textit{recall} e precisão estatisticamente superiores aos do \textit{ruleset} tradicional da ferramenta?
  \end{quote}
  \vspace{0.3cm}
  \textbf{Escopo:} linguagem C, SemGrep como motor de execução, cinco CWEs de memória:
  \begin{center}
  \small CWE-401 (vazamento) \quad CWE-415 (liberação dupla) \quad CWE-416 (uso após liberação)\\
  CWE-457 (variável não inicializada) \quad CWE-476 (ponteiro nulo)
  \end{center}
\end{frame}

% ---------------------------------------------------------------
\section{Ferramentas utilizadas}
% ---------------------------------------------------------------
\begin{frame}{Ferramenta: SemGrep}
  \footnotesize
  \begin{itemize}
    \item Analisador estático de código aberto: regras em YAML, com padrões que imitam a própria sintaxe da linguagem analisada -- ao contrário do CodeQL, que trata o código como um banco de dados relacional.
    \item Converte o código em uma árvore de sintaxe abstrata universal (UAST), com \textit{pattern-matching} uniforme em mais de 30 linguagens.
    \item \textbf{Metavariáveis} (\texttt{\$VAR}) capturam um valor e exigem reuso idêntico, permitindo fluxo de dados; \textbf{reticências} (\texttt{...}) abstraem qualquer trecho irrelevante do padrão.
    \item Modo \textit{taint}: declara fonte, sumidouro (\textit{sink}) e, opcionalmente, desinfetantes (\textit{sanitizers}) de contaminação -- o modo usado neste trabalho.
    \item Mantém um \textit{ruleset} público/comunitário -- é esse conjunto que este trabalho chama de ``regra tradicional''.
  \end{itemize}
\end{frame}

\begin{frame}{Exemplo de regra SemGrep}
  \centering
  \includegraphics[width=0.66\textwidth]{Images/Semgrep.png}
  \vspace{-0.1cm}

  \footnotesize Regra que identifica uma senha embutida como valor padrão de um argumento, em Python; o predicado \texttt{pattern-not} exclui o caso em que a senha está vazia.
  \vspace{-0.05cm}

  {\scriptsize\itshape Fonte: adaptado de Hou et al. (2025) -- mesma figura do TCC (Seção 1.1).\par}
\end{frame}

\begin{frame}{Ferramenta: Juliet Test Suite}
  \begin{itemize}
    \item Coleção de domínio público do NIST (SARD): dezenas de milhares de casos de teste sintéticos em C/C++ e Java, cada um exemplificando uma vulnerabilidade categorizada por CWE.
    \item Cada caso de teste existe em par: uma variante vulnerável e uma variante corrigida do mesmo código.
    \item Papel neste trabalho: (1) fonte de treino/teste da validação cruzada ($k=5$); (2) piso de qualidade do \textit{gate} -- zero falsos positivos exigidos sobre o Juliet.
    \item Também alimenta o modo de exemplos \texttt{pairs} do gerador: os pares vulnerável/seguro do próprio Juliet, mostrados como exemplo ao LLM.
  \end{itemize}
\end{frame}

% ---------------------------------------------------------------
\section{Objetivos e método}
% ---------------------------------------------------------------
\begin{frame}{Objetivos e hipótese}
  \textbf{Objetivo geral:} Um \textit{framework} de síntese automatizada de regras com \textit{recall} e precisão superiores ao \textit{ruleset} tradicional, em pelo menos cinco categorias de CWE.
  \vspace{0.2cm}

  \textbf{Objetivos específicos:}
  \begin{enumerate}
    \item Delimitar os limites de detecção semântica da abordagem gerativa.
    \item Medir o ganho quantitativo de \textit{Recall}, Precisão e F1-Score.
    \item Avaliar a viabilidade de modelos de menor escala em infraestrutura local.
  \end{enumerate}
  \vspace{0.2cm}

  \textbf{Hipótese:} o modelo automatizado detecta vulnerabilidades com taxas estatisticamente superiores às abordagens tradicionais.
\end{frame}

\begin{frame}{Arquitetura implementada}
  \begin{center}
  \begin{tikzpicture}[
      node distance=0.9cm and 1.5cm,
      every node/.style={font=\small},
      box/.style={draw, rounded corners, fill=senacblue!10, align=center, minimum height=1.1cm, minimum width=2.4cm}
    ]
    \node[box] (ger) {Gerador\\(LLM)};
    \node[box, right=of ger] (gate) {\textit{Gate}\\determinístico};
    \node[box, right=of gate] (final) {Conjunto final\\de regras};
    \node[box, below=0.9cm of gate] (cor) {Corretor\\(LLM)};

    \draw[->, thick] (ger) -- node[above,font=\scriptsize]{regra candidata} (gate);
    \draw[->, thick] (gate) -- node[above,font=\scriptsize]{aprovada} (final);
    \draw[->, thick] (gate.south) -- node[right,font=\scriptsize]{reprovada} (cor.north);
    \draw[->, thick] (cor.west) .. controls +(-0.8,0.25) and +(-0.8,-0.25) .. node[left,font=\scriptsize,align=right]{nova\\versão} (gate.west);
  \end{tikzpicture}
  \end{center}
  \vspace{-0.3cm}
  \begin{itemize}
    \footnotesize
    \item \textbf{\textit{Gate}}: valida a sintaxe, executa o SemGrep e mede \textit{recall}, precisão e falsos positivos.
    \item \textbf{Aprovação}: \textit{recall} $\geq$ 50\%, zero falsos positivos no Juliet, $\leq$ 0,1 alertas por KLOC em código real.
    \item Até 6 tentativas de correção -- teto empírico (começou em 3); a repetição persistiu mesmo assim, o que motivou o \texttt{field\_correct} (corrige só o campo errado, não a regra inteira).
    \item Validação cruzada ($k=5$), até 5 sementes por configuração.
  \end{itemize}
\end{frame}

% ---------------------------------------------------------------
\section{Resultados}
% ---------------------------------------------------------------
\begin{frame}{Antes dos números: formato e combinação de modelo}
  \begin{itemize}
    \item \textbf{Formato de saída} (liberdade decrescente do modelo gerador): \texttt{yaml} $>$ \texttt{spec} $>$ \texttt{template} $>$ \texttt{form} $>$ \texttt{formsearch}.
    \item \textbf{Combinação de modelo} (\textit{combo}): define o(s) LLM(s) usados como gerador/corretor -- Combo D usa um modelo de 30B (\textit{Mixture-of-Experts}); Combo C, um modelo de 4B.
  \end{itemize}
  \vspace{0.2cm}
  \textit{\footnotesize Leitura: ``Combo D, \texttt{template}'' = modelo de 30B, formato \texttt{template}.}
\end{frame}

\begin{frame}{Nossas combinações de modelo}
  \footnotesize
  \begin{table}
  \begin{tabular}{cp{9cm}}
  \toprule
  \textbf{Combo} & \textbf{Modelo gerador/corretor} \\
  \midrule
  A & Qwen2.5-Coder-3B em todos os papéis \\
  B & Qwen2.5-Coder-7B \\
  C & Qwen3-4B-Instruct-2507 \\
  D & Qwen3-Coder-30B-A3B (\textit{MoE}, $\sim$3B parâmetros ativos/\textit{token}) \\
  E & Qwen3-Coder-30B-A3B, sem troca de modelos entre papéis \\
  \bottomrule
  \end{tabular}
  \end{table}
  \vspace{-0.1cm}
  \begin{itemize}
    \item Todas usam o mesmo modelo de detecção inicial (Qwen2.5-Coder-3B); variam só o gerador/corretor.
    \item Combo D foi necessária para os melhores resultados; Combo C (4B) não reproduziu os ganhos do \texttt{field\_correct}.
  \end{itemize}
\end{frame}

\begin{frame}{Modo de exemplos: \texttt{findings} (padrão)}
  \begin{itemize}
    \item O \textbf{modo de exemplos} (\texttt{example\_mode}) define o que é mostrado ao gerador antes de escrever a regra -- há quatro valores: \texttt{findings}, \texttt{pairs}, \texttt{graph} e \texttt{none}.
    \item Em \texttt{findings}, o gerador recebe o relatório textual de um \textbf{detector inicial}: um LLM menor (3B) que descreve, em linguagem natural, onde acredita haver o defeito.
    \item Foi o modo padrão, usado nas primeiras rodadas experimentais, antes de qualquer alternativa ser testada.
    \item Limitação observada: a descrição é predominantemente textual -- não aponta as linhas exatas do código envolvidas no defeito.
  \end{itemize}
\end{frame}

\begin{frame}{Modo de exemplos: \texttt{pairs}}
  \begin{itemize}
    \item O gerador recebe, lado a lado, o par de código do próprio Juliet Test Suite: a versão \textbf{vulnerável} e a versão \textbf{corrigida} do mesmo trecho.
    \item É o parâmetro da \textbf{única configuração efetivamente aprovada pelo \textit{gate}} ao longo de todo o trabalho (Combo D, \texttt{template}, CWE-416).
    \item Vantagem sobre \texttt{findings}: mostra ao modelo não só onde está o defeito, mas também o que o diferencia do código seguro equivalente.
  \end{itemize}
  \vfill
  {\tiny\itshape CWE-416: Uso Após Liberação (\textit{Use After Free}).\par}
\end{frame}

\begin{frame}{Modo de exemplos: \texttt{graph}}
  \begin{itemize}
    \item Alimentado por um componente à parte, sem LLM: o analisador de grafos de eventos de memória, que aponta as duas linhas exatas do defeito (liberação e uso).
    \item Efeito medido no CWE-416 (Combo D): \textit{recall} sobe, mas a regra fica mais ruidosa -- mais abrangente, não mais silenciosa.
  \end{itemize}
  \begin{table}
  \small
  \begin{tabular}{lcc}
  \toprule
  \textbf{Métrica} & \textbf{Achados do detector} & \textbf{Achados do grafo} \\
  \midrule
  \textit{Recall} & 71,6\% & 91,6\% \\
  Precisão & 89,4\% & 69,8\% \\
  FP (Juliet) & 9,0 & 29,4 \\
  Alertas/KLOC (mundo real) & 2,53 & 4,76 \\
  \bottomrule
  \end{tabular}
  \end{table}
  \footnotesize Causa: o achado do grafo diz \textit{onde} está o defeito, mas não instrui o modelo a excluir a própria linha de liberação do padrão de uso perigoso.
  \vfill
  {\tiny\itshape CWE-416: Uso Após Liberação (\textit{Use After Free}).\par}
\end{frame}

\begin{frame}{Mecanismos do corretor: \texttt{guard} e \texttt{project\_apis}}
  \begin{itemize}
    \item \textbf{\texttt{guard}}: nunca permite que uma correção substitua a melhor regra já encontrada por uma pior -- protege contra regressões durante a correção.
    \item \textbf{\texttt{project\_apis}}: insere no \textit{prompt} os nomes das funções de memória do próprio projeto-alvo.
    \item Juntos, elevaram o \textit{recall} do CWE-416 -- mas à custa de um salto de ruído; isolar o \texttt{project\_apis} revelou que ele era a fonte do problema (próximo slide).
  \end{itemize}
  \vspace{0.15cm}
  \textit{\footnotesize Leitura: ``\texttt{pairs+guard}'' = modo de exemplos \texttt{pairs} com o mecanismo \texttt{guard} ativo.}
  \vfill
  {\tiny\itshape CWE-416: Uso Após Liberação (\textit{Use After Free}).\par}
\end{frame}

\begin{frame}{Síntese dos resultados por CWE}
  \begin{table}
  \small
  \begin{tabular}{lp{4.4cm}cc}
  \toprule
  \textbf{CWE} & \textbf{Melhor configuração} & \textbf{\textit{Recall}} & \textbf{Aprovação} \\
  \midrule
  416 & Combo D, \texttt{template}, \texttt{pairs+guard} & 89,2--100\% & 2 de 5 sementes \\
  415 & Combo D, \texttt{yaml}, \texttt{graph} & 43,3--76,3\% & 0 de 5 \\
  401 & execução isolada, sem validação cruzada & 79,47\% & -- \\
  476 & Combo D, \texttt{yaml}, \texttt{graph} & 22,9--24,8\% & 0 de 2 \\
  457 & Combo C, \texttt{formsearch} & 0--18,75\% & 0 de 2 \\
  \bottomrule
  \end{tabular}
  \end{table}
  \footnotesize
  \begin{itemize}
    \item Apenas o CWE-416 obteve aprovação consistente do \textit{gate}; o CWE-415 chegou perto, sem jamais ser aprovado.
    \item \textit{Recall} vem do teste isolado do Juliet; aprovação exige também $\leq$ 0,1 alertas/KLOC no git/git (só filtra ruído).
  \end{itemize}
  \vspace{-0.25cm}
  {\tiny\itshape CWE-416: Uso Após Liberação -- CWE-415: Liberação Dupla -- CWE-401: Vazamento -- CWE-476: Ponteiro Nulo -- CWE-457: Var. Não Inicializada.\par}
\end{frame}

\begin{frame}{CWE-416: evolução até a melhor configuração}
  \begin{table}
  \small
  \begin{tabular}{lccc}
  \toprule
  \textbf{Configuração} & \textbf{\textit{Recall}} & \textbf{FP (Juliet)} & \textbf{Sementes aprov.} \\
  \midrule
  \texttt{pairs} & 77,6--97,8\% & 84--320 & 1 de 5 \\
  \texttt{pairs} + \texttt{guard} + \texttt{project\_apis} & 80,3--97,1\% & 619--799 & 0 de 2* \\
  \texttt{pairs} + \texttt{guard} (final) & 89,2--100\% & 12--216 & 2 de 5 \\
  \bottomrule
  \end{tabular}
  \end{table}
  \begin{itemize}
    \footnotesize
    \item *Linha do meio: rodada de diagnóstico isolado, só 2 sementes -- as demais usam 5.
    \item \texttt{project\_apis} (nomes de função do próprio projeto) foi a fonte do ruído -- removê-lo, mantendo apenas \texttt{guard}, deu o melhor resultado do trabalho.
  \end{itemize}
  \vspace{0.2cm}
  {\footnotesize\itshape [Espaço reservado para gráfico -- Images/grafico-fp-cwe416.svg]\par}
  \vfill
  {\tiny\itshape CWE-416: Uso Após Liberação (\textit{Use After Free}).\par}
\end{frame}

\begin{frame}{CWE-415: a lição metodológica das sementes}
  \begin{center}
  \small
  \begin{tabular}{cccccc}
  \toprule
  Semente & 1 & 2 & 3 & 4 & 5 \\
  \midrule
  \textit{Recall} & 63,8\% & 64,8\% & 47,3\% & 43,3\% & 76,3\% \\
  \bottomrule
  \end{tabular}
  \end{center}
  \begin{itemize}
    \item Amplitude de 33 pontos percentuais entre sementes -- maior do que o efeito de qualquer mecanismo testado.
    \item Nunca aprovada pelo \textit{gate}: \textit{recall} $\geq$ 50\% \textbf{e} ruído real $\leq$ 0,1 alertas/KLOC, simultaneamente, nunca aconteceram juntos.
    \item Lição metodológica: avaliar sempre com múltiplas sementes antes de reportar um resultado.
  \end{itemize}
  \vspace{0.1cm}
  {\footnotesize\itshape [Espaço reservado para gráfico -- Images/grafico-recall-cwe415.svg]\par}
  \vfill
  {\tiny\itshape CWE-415: Liberação Dupla (\textit{Double Free}).\par}
\end{frame}

\begin{frame}{A regra sintetizada supera a regra tradicional}
  \begin{table}
  \small
  \begin{tabular}{llccc}
  \toprule
  \textbf{CWE} & \textbf{Abordagem} & \textbf{\textit{Recall}} & \textbf{Precisão} & \textbf{F1} \\
  \midrule
  \multirow{2}{*}{416} & Tradicional (\texttt{r/c}) & 17,4\% & 66,7\% & 0,276 \\
   & Sintetizada & 89,2--100\% & 56,1--92,4\% & 0,72--0,91 \\
  \midrule
  \multirow{2}{*}{415} & Tradicional (\texttt{r/c}) & 28,7\% & 50,0\% & 0,364 \\
   & Sintetizada & 43,3--76,3\% & 41,6--59,7\% & 0,50--0,63 \\
  \bottomrule
  \end{tabular}
  \end{table}
  \begin{itemize}
    \footnotesize
    \item \texttt{r/c}: pacote do registro público do SemGrep -- a ``regra tradicional'' já apresentada.
    \item Tradicional é um valor único (testada sobre todo o \textit{dataset} rotulado); Sintetizada é uma faixa (validação cruzada, 5 sementes).
    \item \textit{Recall} e F1 superiores em \textbf{todas} as cinco sementes, nos dois CWEs, sem exceção.
  \end{itemize}
  \vspace{0.1cm}
  {\footnotesize\itshape [Espaço reservado para gráfico -- Images/grafico-comparacao-tradicional.svg]\par}
  \vfill
  {\tiny\itshape CWE-416: Uso Após Liberação -- CWE-415: Liberação Dupla.\par}
\end{frame}

\begin{frame}{O que não funcionou: um limite da linguagem, não da busca}
  \begin{itemize}
    \small
    \item \textbf{CWE-457}: o motor \textit{Open Source} do SemGrep não inicia rastreamento de contaminação a partir de uma declaração sem valor -- teto de \textit{recall} conhecido ($\approx$31\%), abaixo do piso do \textit{gate}.
    \item \textbf{CWE-401} (vazamento): a condição é um argumento de vivacidade sobre \textbf{todos} os caminhos da função -- não uma relação entre dois pontos fixos.
    \item \textbf{CWE-415} (liberação dupla): fonte e sumidouro de contaminação são o \textbf{mesmo padrão} (\texttt{free(\$P)}); o motor de \textit{taint} casa a ocorrência contra si mesma.
    \item \textbf{CWE-476} (ponteiro nulo): a ausência de teste é uma condição sobre \textbf{caminhos}, não uma forma sintática única a excluir.
  \end{itemize}
  \vspace{0.02cm}
  \footnotesize Confirmado por três métodos independentes -- modelo de 30B; gerador determinístico \textbf{sem LLM}; regra fixa compilada direto do grafo, \textbf{sem busca} -- o mesmo limite se repete nos três.
  \vfill
  {\tiny\itshape CWE-457: Variável Não Inicializada -- CWE-401: Vazamento de Memória -- CWE-415: Liberação Dupla -- CWE-476: Ponteiro Nulo.\par}
\end{frame}

% ---------------------------------------------------------------
\section{Conclusão}
% ---------------------------------------------------------------
\begin{frame}{Conclusão}
  \begin{itemize}
    \item \textbf{Hipótese parcialmente confirmada}: refutada na forma ampla (nem todas as 5 categorias); \textbf{não refutada} nas duas em que havia uma regra tradicional para comparar (CWE-416 e CWE-415).
    \item Objetivos 1 (limites semânticos) e 3 (viabilidade local) \textbf{plenamente alcançados}; Objetivo 2 (comparação quantitativa) \textbf{parcialmente alcançado}.
    \item Contribuição central: distinguir um limite de \textbf{busca} de um limite da \textbf{linguagem-alvo} do SemGrep -- generalizável para além deste trabalho.
    \item Achado complementar: aprovação pelo \textit{gate} não garante alertas verdadeiros em código real -- revisão manual continua necessária.
    \item \textbf{Próximos passos:} estender a validação de negativos do mundo real para além do \textit{git/git} (hoje o único repositório usado no \textit{gate}); explorar motores com sensibilidade a caminho (CodeQL, Joern) para os CWEs sem solução.
  \end{itemize}
\end{frame}

\begin{frame}[plain]
  \centering
  \vspace{1cm}
  \logosenacgrande

  \vspace{0.8cm}
  {\Large\bfseries\color{senacblue}Obrigado!\par}
  \vspace{0.3cm}
  {\normalsize Perguntas?\par}
  \vspace{0.6cm}
  {\footnotesize Caio Xavier da Silva \\ \texttt{[caioxaviercxs@gmail.com]}}
\end{frame}

% ---------------------------------------------------------------
% APÊNDICE (slides de apoio para a arguição da banca -- fora dos 15 min)
% ---------------------------------------------------------------
\appendix
\begin{frame}{Apêndice: slides de apoio}
  Referencial teórico, trabalhos correlatos, formatos de regra,
  combinações de modelo, o analisador de grafos de eventos e as
  contribuições detalhadas (DSR) -- usar apenas se a banca perguntar.
\end{frame}

\begin{frame}{Apêndice -- Trabalhos correlatos}
  \small
  \begin{itemize}
    \item \textbf{KNighter} (Yang et al., 2025): gera \textit{checkers} para o Clang Static Analyser a partir de \textit{commits} de correção; pioneiro, mas 32,2\% dos novos alertas eram falsos positivos.
    \item \textbf{Hou et al. (2025)}: gera/migra/otimiza regras SemGrep via \textit{few-shot learning}; 41 regras válidas contra 37 do KNighter.
    \item \textbf{RepoAudit} (Guo et al., 2025): agente autônomo com memória e validação por \textit{tree-sitter}; precisão de 78,43\%, mas reconsulta o LLM a cada auditoria.
    \item \textbf{Este trabalho}: LLM só na geração da regra -- a execução final não tem custo de inferência.
  \end{itemize}
\end{frame}

\begin{frame}{Apêndice -- O analisador de grafos de eventos}
  \small
  \begin{itemize}
    \item Componente complementar, em \textit{branch} separado, \textbf{sem nenhum LLM}.
    \item Constrói, via \textit{tree-sitter}, um grafo de eventos de memória por função (alocação, liberação, uso, reatribuição) e faz consultas de alcançabilidade.
    \item Nomes de função de alocação/liberação inferidos estatisticamente do próprio projeto -- sem lista manual.
    \item Isolado: 134 de 138 casos de CWE-416 do Juliet, zero falsos positivos; aprovado em \textit{tmux}, PostgreSQL e SQLite; falhou no \textit{vim} (uma função de liberação mal inferida gerou 179 dos 234 alertas).
    \item Usado também como fonte de exemplos para o gerador (\texttt{example\_mode=graph}).
  \end{itemize}
  \vfill
  {\tiny\itshape CWE-416: Uso Após Liberação (\textit{Use After Free}).\par}
\end{frame}

\begin{frame}{Apêndice -- Formatos de regra e mecanismos avaliados}
  \small
  \textbf{Formatos de saída} (liberdade decrescente): \texttt{yaml}, \texttt{spec}, \texttt{template}, \texttt{form}, \texttt{formsearch}.
  \vspace{0.15cm}

  \textbf{\texttt{field\_correct}}: corrige um único campo por tentativa, em vez de reescrever a regra inteira -- ganho real, porém dependente da combinação de CWE, formato e porte do modelo; sem generalização.
  \vspace{0.15cm}

  \textbf{Combinações de modelo}: A (3B) a E (30B, \textit{Mixture-of-Experts}); a combinação D foi necessária para os melhores resultados -- modelos de 4B (combo C) são viáveis localmente, mas com desempenho inferior em mecanismos de correção mais refinados.
\end{frame}

\begin{frame}{Apêndice -- Contribuições (Design Science Research)}
  \small
  \begin{itemize}
    \item \textbf{Constructo}: caracterização do compromisso entre liberdade sintática e eficácia, nos cinco formatos de regra.
    \item \textbf{Método}: caracterização do \texttt{field\_correct} como ganho condicional, não geral.
    \item \textbf{Modelo}: evidência de que a variação entre sementes pode exceder o efeito de qualquer mecanismo experimental.
    \item \textbf{Instanciação}: \textit{pipeline} completo e executável em \textit{hardware} de consumo, com modelos locais de código aberto.
  \end{itemize}
\end{frame}

\end{document}
